\chapter{问题、例子与主定理}\label{sec:main}

\section{从一个小实验开始}
设有七个带编号的位置，数值为
\[
(-3,0,0,0,1,1,1).
\]
每次选择三个位置，将三个数同时替换成它们的平均。
先平均 $(-3,0,0)$，便得到三个 $-1$、三个 $1$ 和一个 $0$。
接着三次选取 $(-1,0,1)$，所有位置都变成零。
这个例子体现两个基本事实：同样的数可以占据多个不同位置；
某次平均产生的零，可以在后面的步骤中继续使用。

我们研究的是：给定有理数配置，是否存在某个\emph{有限}平均序列，使所有位置精确相等。
选取位置的顺序可以依赖输入。不断接近相等还不足以回答问题，因为极限相等与有限步相等是不同的要求。

\section{形式定义与本原中心化}

配置是一个有序列 $x=(x_1,\ldots,x_n)\in\Q^n$。对 $p$ 元位置集 $S$，令
\[
(T_Sx)_i=\begin{cases}p^{-1}\sum_{j\in S}x_j,&i\in S,\\x_i,&i\notin S.\end{cases}
\]
这些映射称为 $p$ 平均。其共同固定的线性泛函是总和，故全等终点只能是 $\bar x\mathbf1$，其中 $\bar x=n^{-1}\sum_i x_i$。


总和不变直接来自所选组的总和等于其均值的 $p$ 倍。因此终点没有自由度，
只可能是原全局均值。把这个均值减去以后，目标便变为零。
再乘一个共同整数可以清除所有分母，最后除去各坐标的最大公因子。
这个过程只改变我们描述输入所用的尺度，不改变哪一串平均操作能够成功。
下一讲将把这种相容性写成正式引理。

若 $x$ 非恒定，选择 $c\in\Q^\times$，使
\[
X=c(x-\bar x\mathbf1)\in\Z^n,\qquad \sum_iX_i=0,\qquad \gcd_iX_i=1.
\]
$X$ 唯一至整体符号，称为 $x$ 的本原中心化代表。定义
\begin{equation}\label{eq:G}
G(x)=G(X)=\gcd_{i<j}|X_i-X_j|.
\end{equation}


\begin{example}[如何算判据]
取五平均、十一位置的配置
\[
X=(1^{10},-10).
\]
它已经零和而且本原。任意两个不同坐标的差都是 $11$，所以 $G=11$。
虽然位置数已达到 $2p+1=11$，它仍然不能归零：模 $11$ 的每个数都等于 $1$，
五平均也保持这个余数。
相比之下，$(-5,0^5,1^5)$ 同样有十一位置，且 $G=1$；
先平均 $(-5,0^4)$，便化成相反值和一个零，从而可解。
\end{example}

\section{主定理究竟断言什么}

\begin{theorem}[素数平均的充要条件]\label{thm:prime}
设 $p$ 为奇素数，$n\ge2p+1$。非恒定配置 $x\in\Q^n$ 可经有限次 $p$ 平均达到全等，当且仅当 $G(x)=p^e$，其中 $e\ge0$。对每个奇素数 $p$，在 $n=2p$ 存在 $G=1$ 的不可达配置。
\end{theorem}

\begin{corollary}[最优最终阈值]\label{cor:threshold}
记 $N(p)$ 为使上述充要条件对所有 $n\ge N(p)$ 成立的最小整数。则奇素数满足 $N(p)=2p+1$；二平均满足 $N(2)=4$。
\end{corollary}


这里允许 $e=0$，所以 $G=1$ 属于可解条件。
“最终阈值”要求\emph{以后的每一个}位置数都满足判据。
例如在 $n=p$ 时，一步就能把全部位置平均；这并不影响 $2p+1$ 是最优最终阈值。
阈值以下既可能存在可解配置，也可能存在满足 $G=1$ 却不可解的配置。

\section{证明路线与范围划分}
设 $p\ge5$。下表先说明每一部分在证明中承担什么工作，具体构造随后逐讲给出。
\begin{center}
\begin{tabular}{p{.26\textwidth}p{.61\textwidth}}
\toprule 范围或任务 & 主要方法\\\midrule
必要条件、$n=2p$ 的反例 & 共同余数与最低素数赋值\\
$n\ge4p-2+d_p(n)$ & 保留两个等值大块的整数能量下降\\
$n=3p$ & 三个等大块上的整数剪切\\
$2p<n<4p,\ n\ne3p$ & 一般输入压缩、二维返回、有限同余输送\\
其余 $n\ge4p$ & 加权三块核心及同余输送\\
$p=3,5,7$ 的必要小情形 & 显式平均词、下降及有限分类\\
\bottomrule
\end{tabular}
\end{center}
论证中最重要的先后关系是：\emph{先建立含零收尾，再把临界带送到含零状态。}
这样，后半部分每次“造出零和组”之后，都有已经证明的下一步可用。

本讲义采用如下约定。上标 $a^m$ 表示 $m$ 个取值为 $a$ 的位置；矩阵作用于列向量，复合从右至左。位置始终带标签；交换两个等大块的名称仅改变参数表示。固定平均词与整体仿射变换交换，因而可以在分析中选取本原代表。下文的矩阵返回均覆盖所讨论配置的全部位置。

充分性证明的依赖顺序如下。引理\ref{lem:integer-range}独立处理 $n\ge4p-2+d_p(n)$；命题\ref{prop:threep}独立处理 $n=3p$。二者推出引理\ref{lem:zero-tail}的含零收尾。临界带的入口由命题\ref{prop:entry}给出，核内完成使用定理\ref{thm:terminal}以及命题\ref{prop:odd}, \ref{prop:r2}, \ref{prop:even}, \ref{prop:upper}。不满足整数下降容量的宽维数由命题\ref{prop:wide}补齐。最后在第\ref{sec:completion}讲合并各范围。

矩阵部分使用两个已发表结果：Morris 的有限指数定理\cite[Theorem 6.1(2)]{morris}，以及有理 $S$ 整数上 $\SL_2$ 的强同余子群性质\cite{serre,prasad}。它们的精确代入见定理\ref{thm:arithmetic-external}。有限证书的输入全集及检验条件见附录\ref{app:certificates}。

